% TEMPLATE-MILC-v3.tex - plantilla maestra UNICA de las guias MILC v3.
%
% La usan las cuatro familias de guias, cada una con su driver:
%   · Grados 6-11          -> scripts/build-guias-g11.py
%   · Bebras               -> scripts/build-guias-bebras.py
%   · Semillero            -> scripts/build-guias-semillero.py
%   · Territorio interior  -> scripts/build-guias-territorio-interior.py
%
% Antes habia tres copias casi identicas (~400 lineas iguales, 6 distintas):
% cada mejora habia que aplicarla tres veces, y en la practica no se hacia
% (el motor de recursos vivia solo en dos de ellas). Lo que varia entre
% familias esta parametrizado en 6 placeholders que cada driver llena
% (se nombran aqui SIN su sintaxis real: el builder reemplaza por texto plano,
% tambien dentro de los comentarios, y un valor multilinea rompe el preambulo):
%
%   HEADER_IZQ        encabezado izquierdo de pagina
%   HEADER_DER        encabezado derecho de pagina
%   PORTADA_ETIQUETA  rotulo sobre el numero grande (GRADO/MOMENTO/MODULO)
%   PORTADA_SERIE     linea de serie de la portada
%   PORTADA_ID        identificador de la guia en la portada
%   PORTADA_CIRCULO   el circulito de grado (°), o vacio si no aplica
%   PDF_TITULO        titulo en texto plano (sin \textbf ni comillas TeX) para
%                     los metadatos del PDF (/Title); lo produce lib_xelatex.texto_pdf
%
% Composicion (auditoria 2026-09-03): las bandas de estacion piden espacio
% (\Needspace*) y prohiben el salto justo despues (\nopagebreak), asi no quedan
% huerfanas al pie; las cajas largas (softbox, drawbox, infoband, puentebox) son
% `breakable`, asi una caja de media pagina ya no arrastra todo a la siguiente
% dejando la anterior al 40 %. Todo texto sobre mostaza o verde va en milcNegro
% (blanco sobre #E5B400 da 1,93:1 y sobre #58A923 2,95:1; ninguno pasa WCAG AA).
% El resto de claves las documenta content/guias/_SCHEMA.md. El script valida
% que no queden placeholders sin reemplazar antes de escribir el archivo final.

% Etiquetado PDF (latex-lab): /Lang, estructura y texto alternativo de las
% figuras, para lectores de pantalla. Debe ir ANTES de \documentclass.
\DocumentMetadata{lang=es-CO,tagging=on}
\documentclass[11pt,letterpaper]{article}
% headheight=24pt: la cabecera derecha lleva dos lineas (identidad de la guia
% y estacion actual). headsep se acorta en la misma medida para que el bloque
% de texto quede exactamente donde estaba con headheight=12pt.
\usepackage[margin=2.0cm,headheight=24pt,headsep=13pt]{geometry}
\usepackage[table]{xcolor}
\usepackage{fontspec}
\usepackage[spanish]{babel}
\usepackage{tikz}
% Librerías TikZ para los diagramas de las guías (bloque `recursos:`):
%  · babel  → babel en español vuelve activos caracteres como «>», y TikZ los usa
%             en sus opciones (>=stealth, ->). Sin esta librería, cualquier
%             diagrama con flechas revienta con «\language@active@arg> has an
%             extra }». Debe ir ANTES de que se dibuje nada.
%  · positioning        → sintaxis «below left=2pt and 3pt».
%  · decorations.pathreplacing → llaves ({brace}) para señalar rangos.
\usetikzlibrary{babel,positioning,decorations.pathreplacing}
\usepackage{graphicx}
% Circuitos: el trabajo de electrónica se hace hoy con micro:bit y sus sensores
% integrados (MakeCode, sin cableado), así que no se necesitan esquemáticos.
% Si algún día entran montajes con protoboard, basta descomentar esta línea:
% los diagramas .tex/.tikz declarados en `recursos:` ya se incrustan solos.
% \usepackage{circuitikz}
\usepackage[most]{tcolorbox}
\usepackage{tabularx}
\usepackage{array}
\usepackage{fancyhdr}
\usepackage{titlesec}
\usepackage[document]{ragged2e}  % bandera a la derecha en todo el cuerpo
\usepackage{enumitem}
\usepackage{microtype}
\usepackage{etoolbox}
\usepackage{needspace}
% hyperref al final del preambulo: metadatos del PDF (titulo, autor, idioma,
% familia), marcadores por estacion (\pdfbookmark) y enlaces sin recuadro.
\usepackage[hidelinks,
  pdftitle={Cosecha del periodo 2 — mi primer programa Scratch completo},
  pdfauthor={PhD. Álvaro Cárdenas Orozco},
  pdfsubject={Tecnología e Informática · Grado 7 · Guía 20 · MILC v3},
  pdflang={es-CO},
  bookmarksopen=true,bookmarksnumbered=false]{hyperref}

% Helvetica Neue no trae la flecha U+2192, asi que cada «→» escrito literalmente
% en el YAML se compilaba a NADA, en silencio (462 flechas en 61 guias: rutas del
% tipo «paleta → tipografia → lockup» salian rotas). Mapeamos el caracter a su
% version matematica, que si tiene glifo. Asi el autor puede seguir escribiendo
% «→» comodamente en el YAML.
\usepackage{newunicodechar}
\newunicodechar{→}{\ensuremath{\rightarrow}}
% Mismo problema con los simbolos de marca y vinetas: el «✓» de «una palomita ✓»
% salia como cuadrito vacio en el PDF de todas las guias que lo usaban.
\usepackage{pifont}
\newunicodechar{✓}{\ding{51}}
\newunicodechar{✗}{\ding{55}}
\newunicodechar{✔}{\ding{52}}
\newunicodechar{•}{\textbullet}
\newunicodechar{≥}{\ensuremath{\geq}}
\newunicodechar{≤}{\ensuremath{\leq}}

\setmainfont{Helvetica Neue}
\setsansfont{Helvetica Neue}
\setlength{\parindent}{0pt}
\setlength{\parskip}{6pt}
% Sin particion de palabras: en 6.º-7.º una palabra cortada por guion al final
% de linea («Comuni-cacion») frena la lectura mucho mas que una linea corta.
\hyphenpenalty=10000
\exhyphenpenalty=10000
\pretolerance=2000
\tolerance=3000
\setlength{\emergencystretch}{3em}
\linespread{1.10}
\renewcommand{\arraystretch}{1.25}
\setlist{leftmargin=5mm,itemsep=1mm,topsep=1mm}
\newcolumntype{Y}{>{\RaggedRight\arraybackslash}X}

\definecolor{milcMagenta}{HTML}{E6007E}
\definecolor{milcVino}{HTML}{5A0038}
\definecolor{milcTurquesa}{HTML}{0093A5}
\definecolor{milcVerde}{HTML}{58A923}
\definecolor{milcMostaza}{HTML}{E5B400}
\definecolor{milcOcre}{HTML}{B86600}
\definecolor{milcNegro}{HTML}{241816}
\definecolor{milcGris}{HTML}{F7F7F4}
\definecolor{milcLinea}{HTML}{E8E2DD}
\definecolor{milcGradePrimary}{HTML}{84CC16}
\definecolor{milcGradeDark}{HTML}{4D7C0F}
\definecolor{milcGradeSoft}{HTML}{ECFCCB}
\definecolor{milcGradeAccent}{HTML}{CFE7FF}
\definecolor{milcGradeText}{HTML}{FFFFFF}
\color{milcNegro}

\pagestyle{fancy}
\fancyhf{}
% Cabecera de dos lineas: identidad de la guia arriba y, a la derecha y en
% gris, la estacion en curso (\markright desde \milcestacion). Antes de la
% primera estacion la segunda linea va vacia; el \strut mantiene la altura
% para que la linea de arriba no baile entre paginas.
\lhead{\small Tecnología e Informática · Grado 7\\[-1pt]{\footnotesize\strut}}
\rhead{\small Guía 20 · MILC v3\\[-1pt]{\footnotesize\strut\color{milcNegro!65}\rightmark}}
\cfoot{\small\thepage}
\renewcommand{\headrulewidth}{0pt}

% Sobre mostaza (#E5B400) y verde (#58A923) el texto va en milcNegro; sobre el
% resto de la paleta, en blanco. Se usa dentro de fonttitle/fontupper, donde un
% \color posterior manda sobre coltitle.
\newcommand{\milctintasobre}[1]{%
  \ifboolexpr{test{\ifstrequal{#1}{milcMostaza}} or test{\ifstrequal{#1}{milcVerde}}}%
    {\color{milcNegro}}{\color{white}}}

\newtcolorbox{titlebox}[1]{
  enhanced,colback=#1,colframe=#1,arc=7mm,boxrule=0pt,
  left=7mm,right=7mm,top=3mm,bottom=3mm,
  before skip=8pt,after skip=8pt,
  fontupper=\sffamily\bfseries\Large\color{white}
}

\newtcolorbox{softbox}[3]{
  enhanced,breakable,pad at break=3mm,
  colback=#2,colframe=#1,borderline west={4pt}{0pt}{#1},
  arc=2mm,boxrule=.4pt,left=4mm,right=4mm,top=3mm,bottom=3mm,
  before skip=6pt,after skip=8pt,title={#3},coltitle=white,
  colbacktitle=#1,fonttitle=\sffamily\bfseries\milctintasobre{#1},fontupper=\RaggedRight,
  attach boxed title to top left={xshift=3mm,yshift=-2mm},
  boxed title style={arc=2mm, boxrule=0pt}
}

\newtcolorbox{drawbox}[2]{
  enhanced,breakable,pad at break=4mm,
  colback=white,colframe=#1,arc=4mm,boxrule=1pt,
  left=5mm,right=5mm,top=4mm,bottom=4mm,before skip=8pt,after skip=8pt,
  title={#2},coltitle=white,colbacktitle=#1,fonttitle=\sffamily\bfseries\milctintasobre{#1},
  fontupper=\RaggedRight,
  attach boxed title to top left={xshift=4mm,yshift=-2mm},
  boxed title style={arc=3mm, boxrule=0pt}
}

\newtcolorbox{infoband}[1]{
  enhanced,breakable,pad at break=2mm,colback=#1!8,colframe=#1!50,arc=2mm,boxrule=.5pt,
  left=4mm,right=4mm,top=2mm,bottom=2mm,before skip=4pt,after skip=4pt,
  fontupper=\small\RaggedRight
}

% Puente narrativo entre secciones (contrato editorial Conéctate).
% Da continuidad: "Acabas de X. Ahora vas a Y porque Z."
% Visualmente: banda izquierda en color del grado, texto en itálica pequeña.
\newtcolorbox{puentebox}{
  enhanced,breakable,pad at break=2mm,colback=milcGradeSoft,colframe=milcGradeDark,
  borderline west={3pt}{0pt}{milcGradeDark},
  arc=0mm,boxrule=0pt,left=5mm,right=4mm,top=2.5mm,bottom=2.5mm,
  before skip=4pt,after skip=8pt,
  fontupper=\itshape\small\color{milcNegro}\RaggedRight
}
% La flecha va en modo matematico a proposito: Helvetica Neue NO tiene el glifo
% U+2192, asi que \textrightarrow se compilaba a NADA («Missing character: There
% is no → in font Helvetica Neue») y el puente salia sin flecha. En modo
% matematico la flecha viene de la fuente matematica y si se dibuja.
\newcommand{\puente}[1]{%
  \begin{puentebox}%
    \textcolor{milcGradeDark}{$\rightarrow$}\hspace{2mm}#1%
  \end{puentebox}%
}

\newcommand{\linead}[1]{\noindent\color{milcLinea}\rule{#1}{0.5pt}\color{milcNegro}}
\newcommand{\respuesta}[1]{\par\vspace{2mm}\foreach \n in {1,...,#1}{\noindent\color{milcLinea}\rule{\linewidth}{0.5pt}\par\vspace{5mm}}\color{milcNegro}}
\newcommand{\checkbox}{\textcolor{milcGradeDark}{\fbox{\phantom{x}}}\hspace{5pt}}

% ─── Listas aireadas para las guías socioemocionales ────────────────────
% Componer las verificaciones como «\checkbox texto\\» deja el texto que salta
% de línea debajo del cuadrito y apelmaza el bloque. Estos entornos alinean la
% continuación y airean los ítems. Los usa Territorio Interior; cualquier
% familia puede hacerlo.
\newlist{checklista}{itemize}{1}
\setlist[checklista]{
  label=\textcolor{milcGradeDark}{\fbox{\phantom{x}}},
  leftmargin=9mm, labelsep=3.2mm, itemsep=2.4mm,
  topsep=2.5mm, parsep=0pt, partopsep=0pt, before=\RaggedRight
}
\newlist{pasoslista}{enumerate}{1}
\setlist[pasoslista]{
  label=\textcolor{milcGradeDark}{\sffamily\bfseries\arabic*.},
  leftmargin=8mm, labelsep=3mm, itemsep=2.8mm,
  topsep=1mm, parsep=0pt, partopsep=0pt, before=\RaggedRight
}

\newcommand{\phasecard}[4]{
\begin{tcolorbox}[enhanced,colback=#3,colframe=#2,arc=3mm,boxrule=.5pt,height=3.05cm,valign=center,halign=center,left=1.5mm,right=1.5mm,top=2mm,bottom=2mm]
{\sffamily\bfseries\fontsize{22}{24}\selectfont\textcolor{#2}{#1}}\par
{\sffamily\bfseries\small #4}
\end{tcolorbox}
}

% ── Piezas del rediseno 2026 (legibilidad 6.º-11.º) ───────────────────
% Banda de estacion: reemplaza el titlebox macizo. Titulo a la izquierda,
% dato practico (tiempo/modalidad) a la derecha, en una sola linea.
% Nunca huerfana: pide 9 lineas libres (o salta de pagina rellenando con
% \vfill) y prohibe el salto justo despues de la banda. Nueve y no seis:
% la caja partible que sigue necesita ~80pt (titulo adosado, relleno y dos
% lineas) para arrancar en la misma pagina; con menos, tcolorbox la manda
% entera a la siguiente y la banda se queda sola. El penalty 10000 va
% en `after`, ANTES del espacio posterior: un `after skip` (aunque sea 0pt)
% es un \vskip tras la caja, y ese glue es un punto de corte legal que un
% \nopagebreak escrito despues de \end{tcolorbox} ya no cierra. Cada banda
% deja marcador PDF y actualiza la cabecera derecha.
\newcounter{milcestacion}
\newcommand{\milcestacion}[2]{%
\Needspace*{9\baselineskip}%
\stepcounter{milcestacion}%
\pdfbookmark[1]{#2}{milcestacion\themilcestacion}%
\markright{#2}%
\begin{tcolorbox}[enhanced,colback=#1,colframe=#1,arc=2mm,boxrule=0pt,
  left=5mm,right=5mm,top=2.6mm,bottom=2.6mm,before skip=11pt,
  after={\par\nopagebreak\vskip7pt}]
{\sffamily\bfseries\fontsize{15}{18}\selectfont\milctintasobre{#1}#2}
\end{tcolorbox}}

% Tarjeta de paso numerada: sustituye las tablas «Pilar / Aplicacion», que
% eran formato de consulta docente, no de estudiante. El numero va en un
% circulo sobre el borde para que la secuencia se lea de un vistazo.
\newcommand{\milcpaso}[3]{%
\Needspace{4\baselineskip}%
\begin{tcolorbox}[enhanced,colback=white,colframe=milcLinea,arc=1.5mm,boxrule=.9pt,
  left=14mm,right=4mm,top=3mm,bottom=3mm,before skip=4pt,after skip=4pt,
  fontupper=\RaggedRight,
  overlay={\node[anchor=north west,circle,fill=#2,text=white,inner sep=0pt,
    minimum size=7.4mm,font=\sffamily\bfseries\small]
    at ([xshift=3.6mm,yshift=-3.2mm]frame.north west) {#1};}]
#3
\end{tcolorbox}}

% Casilla marcable de tamano util con lapiz (la anterior era \fbox{\phantom{x}},
% de alto variable segun el texto que la seguia).
\newcommand{\milcmarca}{\framebox[3.8mm]{\rule{0pt}{3.4mm}}}

% Fila marcable de lista suelta (checks de la estacion 1).
\newcommand{\milccheckrow}[1]{%
\par\vspace{1.8mm}\noindent
\makebox[6.4mm][l]{\raisebox{-.55mm}{\textcolor{milcNegro}{\milcmarca}}}%
\parbox[t]{\dimexpr\linewidth-6.4mm\relax}{\RaggedRight #1}\par}

% Celda marcable dentro de la rubrica (tabla de autoevaluacion). Misma
% sangria francesa, pero pensada para vivir dentro de una columna X.
\newcommand{\milcopcion}[1]{%
\makebox[5.2mm][l]{\raisebox{-.55mm}{\textcolor{milcNegro}{\milcmarca}}}%
\parbox[t]{\dimexpr\linewidth-5.2mm\relax}{\RaggedRight #1}}

% ── Recursos visuales de la guía ──────────────────────────────────────
% Declarados en el bloque `recursos:` del YAML y emitidos por el builder.
% \guiaFigura   → imagen rasterizada o PDF (foto, carta celeste, montaje).
% \guiaDiagrama → fuente TikZ/circuitikz (.tex) incrustada como vector:
%                 es la vía para circuitos y esquemáticos editables.
% [#1] = texto alternativo (alt) · #2 = ruta relativa al .tex · #3 = pie
% El alt llega al PDF etiquetado (/Alt) y lo lee el lector de pantalla.
\newcommand{\guiaFigura}[3][]{%
\begin{center}
\includegraphics[width=0.88\linewidth,height=0.38\textheight,keepaspectratio,alt={#1}]{#2}\\[1.5mm]
{\footnotesize\itshape\textcolor{milcNegro!75}{#3}}
\end{center}
}

\newcommand{\guiaDiagrama}[2]{%
\begin{center}
\input{#1}\\[1.5mm]
{\footnotesize\itshape\textcolor{milcNegro!75}{#2}}
\end{center}
}

\begin{document}
\thispagestyle{empty}

% ── Portada ───────────────────────────────────────────────────────────
% Antes: un rectangulo de color a pagina completa. Se veia bien en pantalla,
% pero en la fotocopiadora del colegio se comia el toner y salia gris sucio.
% Ahora la banda ocupa el tercio superior y el resto va en blanco.
\begin{tikzpicture}[remember picture,overlay]
  \path[fill=milcGradePrimary]
    (current page.north west) rectangle ([yshift=-10.6cm]current page.north east);

  \node[anchor=north west,text=milcGradeText] at ([xshift=2.0cm,yshift=-1.55cm]current page.north west)
    {\sffamily\bfseries\fontsize{12}{14}\selectfont GRADO};
  \node[anchor=north west,text=milcGradeText,opacity=.85] at ([xshift=2.0cm,yshift=-2.35cm]current page.north west)
    {\sffamily\bfseries\fontsize{8.5}{10}\selectfont SERIE GUÍAS · TECNOLOGÍA E INFORMÁTICA};
  \node[anchor=north east,text=milcGradeText,opacity=.85,align=right] at ([xshift=-2.0cm,yshift=-1.55cm]current page.north east)
    {\sffamily\bfseries\fontsize{9}{12}\selectfont I.E. Sor María Juliana\\Cartago, Valle del Cauca};

  \node[anchor=west,text=milcGradeText] (gradonum) at ([xshift=1.85cm,yshift=-7.1cm]current page.north west)
    {\sffamily\bfseries\fontsize{112}{112}\selectfont 7};
  % El circulito de grado (°) solo aplica a las guias de grado («11°»); en
  % Bebras y Semillero el numero grande es un momento/modulo, no un grado.
  % Se posiciona relativo al nodo `gradonum`, no en coordenadas absolutas.
  \draw[milcGradeText,line width=1.6pt]
    ([xshift=2.0mm,yshift=-4.5mm]gradonum.north east) circle (.15cm);

  \node[anchor=west,text=milcGradeText,text width=11.4cm,align=left] at ([xshift=1.5cm,yshift=0.5cm]gradonum.east)
    {
     {\sffamily\bfseries\fontsize{9}{11}\selectfont\MakeUppercase{Período 2 · Algoritmia y pensamiento computacional}}\\[3mm]
     {\sffamily\bfseries\fontsize{21}{25}\selectfont\setlength{\baselineskip}{27pt}Cosecha del\\[4pt]periodo 2 ---\\[4pt]mi primer\\[4pt]programa\\[4pt]completo}\\[3.5mm]
     {\sffamily\bfseries\fontsize{15}{18}\selectfont Guía 20-7-TIC}
    };
\end{tikzpicture}

\vspace*{9.4cm}
\vfill

{\sffamily\bfseries\fontsize{8}{10}\selectfont\textcolor{milcGradeDark}{LO QUE TE LLEVAS HOY}}

\begin{tcolorbox}[enhanced,colback=white,colframe=milcNegro,arc=2mm,boxrule=1.6pt,
  left=5mm,right=5mm,top=4mm,bottom=4mm,before skip=3pt,after skip=12pt,
  fontupper=\RaggedRight\fontsize{11}{15.5}\selectfont]
Un \textbf{programa completo en Scratch} que integra TODO lo aprendido en P2. \textbf{Entregables:} \textbf{(a) programa funcional} con: variables (mínimo 2), condicionales (mínimo 2), bucle (mínimo 1), interactividad (preguntar al usuario), múltiples escenas o sprites; \textbf{(b) diagrama de flujo} del programa en el cuaderno; \textbf{(c) documentación} en el cuaderno: descripción del programa, lista de variables, capturas de cada escena, las 4 técnicas del pensamiento computacional aplicadas (S3); \textbf{(d) sustentación} oral de 5 minutos: tu equipo (o tú solo) muestra el programa y explica decisiones de diseño.
\end{tcolorbox}

\begin{tcolorbox}[enhanced,colback=milcGris,colframe=milcGris,arc=2mm,boxrule=0pt,
  left=5mm,right=5mm,top=3.5mm,bottom=3.5mm,after skip=12pt]
{\sffamily\bfseries\fontsize{8}{10}\selectfont\textcolor{milcNegro!55}{ESTA GUÍA ES DE}}\\[3mm]
\begin{tabularx}{\linewidth}{X p{3.2cm} p{3.2cm}}
\linead{\linewidth} & \linead{3.0cm} & \linead{3.0cm}\\[-1mm]
{\fontsize{8.5}{10}\selectfont\textcolor{milcNegro!55}{Nombre completo}} &
{\fontsize{8.5}{10}\selectfont\textcolor{milcNegro!55}{Grupo}} &
{\fontsize{8.5}{10}\selectfont\textcolor{milcNegro!55}{Fecha}}\\
\end{tabularx}
\end{tcolorbox}

\vfill

\noindent\textcolor{milcLinea}{\rule{\linewidth}{1pt}}\\[2mm]
{\sffamily\bfseries\fontsize{9.5}{12}\selectfont Autor: Dr. Álvaro Cárdenas Orozco}\\[1mm]
{\sffamily\fontsize{8.5}{11}\selectfont\textcolor{milcNegro!55}{Metodología MILC · apertura ancestral · pensamiento computacional · triángulo Dussel–estoicismo–Floridi}}

\newpage

\section*{Apertura: Cosecha del periodo 2 --- mi primer programa Scratch completo}

\begin{softbox}{milcGradeDark}{milcGradeSoft}{Saber ancestral}
\textbf{En el Valle del Cauca, cuando una tejedora terminaba una mochila wayuu o un canasto, no se quedaba con él para ella sola.} Doña Sofía y otras tejedoras de la vereda La Plata de Cartago \textbf{mostraban su obra}: \emph{la llevaban a la plaza del pueblo, la enseñaban a los vecinos, contaban qué patrón usaron, cuánto tiempo les tomó, qué cambiarían para la próxima}. La obra terminada se \textbf{sustentaba} ante la comunidad. \emph{Eso no era vanidad}: era validación cultural, era transmisión de saberes, era oportunidad de mejorar. Si alguien notaba un error, lo decía con respeto. Si alguien alababa una innovación, otros aprendían. \textbf{La cosecha no se quedaba en privado}. Hoy tú cierras el periodo 2 igual: \emph{programaste varios miniproyectos en Scratch durante 9 sesiones; hoy presentas tu obra completa, la sustentás, y la dejas como muestra de lo que ahora sabes hacer}.
\end{softbox}

\begin{softbox}{milcTurquesa}{milcGris}{Saber contemporáneo}
Esta es la \textbf{cosecha del periodo 2}. No es contenido nuevo: es \emph{integración} de todo lo aprendido. Tu programa debe combinar: \textbf{algoritmo} (S2 estructura clara con entrada-proceso-salida); \textbf{pensamiento computacional} (S3 las 4 técnicas aplicadas); \textbf{diagrama de flujo} (S4 representación visual); \textbf{variables} (S5 mínimo 2); \textbf{condicionales} (S6 mínimo 2); \textbf{bucles} (S7 mínimo 1); \textbf{Scratch} (S8 conocer la interfaz y bloques); \textbf{narrativa interactiva} (S9 si quieres, preguntar al usuario). \textbf{Tema libre} pero debe ser concreto y demostrable. Sugerencias: \emph{(1) Juego sencillo} (atrapar manzanas, perseguir un sprite, adivinar números); \emph{(2) Historia interactiva extendida} (continuación de la S9 con más decisiones); \emph{(3) Simulador} (el clima del día, el horario semanal); \emph{(4) Animación contadora} (cuenta regresiva con eventos). \textbf{La meta:} demostrar que pasaste de no saber programar a tener un programa funcional propio.
\end{softbox}

\begin{softbox}{milcMostaza}{milcGris}{Pregunta-puente}
\textit{Hace 10 sesiones \emph{no sabías qué era un algoritmo}. Hoy programas en Scratch un programa con variables, condicionales y bucles. \textbf{¿Cómo te sientes con ese crecimiento}?}
\end{softbox}

\begin{softbox}{milcVerde}{milcGris}{Saber y saber hacer}
Al terminar la clase: (1) podrás \textbf{integrar} todos los componentes aprendidos en un solo programa; (2) sabrás \textbf{aplicar} las 4 técnicas del pensamiento computacional en tu diseño; (3) podrás \textbf{evaluar} tu programa con rúbrica; (4) habrás \textbf{creado} y sustentado tu primer programa Scratch completo.
\end{softbox}



\small
\begin{tabularx}{\textwidth}{>{\bfseries}p{3.0cm}Y}
\rowcolor{milcGradePrimary!12}Autor & Dr. Álvaro Cárdenas Orozco\\
DBA & Construye algoritmos y diagramas de flujo para resolver problemas paso a paso (MEN, T\&I 7°)\\
Referentes & Saber ancestral del tejedor · Pensamiento computacional · Scratch\\
Producto final & Un \textbf{programa completo en Scratch} que integra TODO lo aprendido en P2. \textbf{Entregables:} \textbf{(a) programa funcional} con: variables (mínimo 2), condicionales (mínimo 2), bucle (mínimo 1), interactividad (preguntar al usuario), múltiples escenas o sprites; \textbf{(b) diagrama de flujo} del programa en el cuaderno; \textbf{(c) documentación} en el cuaderno: descripción del programa, lista de variables, capturas de cada escena, las 4 técnicas del pensamiento computacional aplicadas (S3); \textbf{(d) sustentación} oral de 5 minutos: tu equipo (o tú solo) muestra el programa y explica decisiones de diseño.\\
\end{tabularx}
\normalsize

\puente{Hoy haces 4 cosas: planeas el programa con las 4 técnicas, lo programas en Scratch, documentas, sustentas.}

\milcestacion{milcGradeDark}{Ruta MILC: así vamos a aprender}
\begin{tabularx}{\textwidth}{>{\centering\arraybackslash}X>{\centering\arraybackslash}X>{\centering\arraybackslash}X>{\centering\arraybackslash}X}
\phasecard{1}{milcVerde}{milcGris}{Escucha\\[-1mm]\small Observo, escucho y pregunto.} &
\phasecard{2}{milcTurquesa}{milcGris}{Entender\\[-1mm]\small Sistematizo y ordeno.} &
\phasecard{3}{milcMagenta}{milcGris}{Hacer\\[-1mm]\small Construyo y aplico.} &
\phasecard{4}{milcVino}{milcGris}{Cierre\\[-1mm]\small Evalúo y reflexiono.}
\end{tabularx}


\puente{Empezamos definiendo el tema y aplicando las 4 técnicas.}

\milcestacion{milcVerde}{Estación 1 de 3 · Escucha}

\begin{softbox}{milcVerde}{milcGris}{Escena para escuchar}
\textbf{Actividad 1 · ANALIZA --- Define tu programa con las 4 técnicas} (12 min · individual). En el cuaderno responde, aplicando las 4 técnicas del pensamiento computacional (S3): \emph{(1) ¿Cuál es tu programa? (1 frase corta).} \emph{(2) DESCOMPOSICIÓN: ¿En qué 4-5 partes se divide? (lista).} \emph{(3) PATRONES: ¿Qué partes son similares entre sí? (que puedan reusar la misma estructura).} \emph{(4) ABSTRACCIÓN: ¿Qué detalles dejas afuera por ahora?}. \emph{(5) ALGORITMO: ¿Cuál es el orden general (entrada-proceso-salida)?}.
\end{softbox}

{\sffamily\bfseries\fontsize{8}{10}\selectfont\textcolor{milcNegro!55}{MARCA CUANDO LO HAYAS HECHO}}
\milccheckrow{Definí el tema concreto de mi programa.}
\milccheckrow{Descompuse en 4-5 partes.}
\milccheckrow{Identifiqué patrones, abstracciones, orden.}

\begin{infoband}{milcVerde}


\textbf{Lo que va al cuaderno (Actividad 1 · ANALIZA).} Título: «Actividad 1 --- Plan del programa con 4 técnicas». \textbf{Formato:} 5 respuestas estructuradas. \textbf{Extensión:} 1/2 página. \textbf{Sabes que terminaste cuando} tu plan está claro antes de tocar Scratch.
\end{infoband}

\puente{Después conoces la rúbrica de 10 criterios + estructura del entregable.}

\milcestacion{milcTurquesa}{Estación 2 de 3 · Entender}

\begin{softbox}{milcTurquesa}{milcGris}{Lo que vas a entender}
Conoces la rúbrica con la que el profe evaluará tu programa. \textbf{Conocerla antes te permite producir según los criterios}.
\end{softbox}

\milcpaso{1}{milcTurquesa}{\textbf{Los 10 criterios de la rúbrica (cada uno 1-5 puntos, total 50).} \textbf{(1) Plan inicial} con las 4 técnicas aplicadas (S3). \textbf{(2) Variables} mínimo 2, bien nombradas según las 5 reglas (S5). \textbf{(3) Condicionales} mínimo 2, con estructura si-entonces-si no (S6). \textbf{(4) Bucle} mínimo 1, no infinito (S7). \textbf{(5) Interactividad} (preguntar al usuario y usar la respuesta). \textbf{(6) Diagrama de flujo} del programa con símbolos correctos (S4). \textbf{(7) Documentación} en cuaderno (descripción + lista de variables + capturas). \textbf{(8) Funcionamiento real} (el programa se ejecuta sin errores). \textbf{(9) Sustentación oral} clara de 5 minutos. \textbf{(10) Personalización} (tema propio, no copia, sprite y fondos personalizados).}
\milcpaso{2}{milcTurquesa}{\textbf{Estructura de tu entregable.} (a) \textbf{Programa en Scratch} guardado en \texttt{scratch.mit.edu} con nombre \texttt{2026-cosecha-p2-[tu-nombre]}. (b) \textbf{Diagrama de flujo} dibujado en el cuaderno aplicando S4 (óvalos, rectángulos, rombos, paralelogramos, flechas). (c) \textbf{Documentación} con: descripción de 1 párrafo, lista de variables con tipo, capturas de 3 momentos del programa, descripción de cómo aplicaste las 4 técnicas. (d) \textbf{Sustentación oral}: 5 minutos donde muestras el programa funcionando, explicas tus decisiones de diseño, comentas qué te costó más.}
\milcpaso{3}{milcTurquesa}{\textbf{3 cosas que diferencian un programa básico de uno bueno.} (1) \emph{Personalización}: tu programa no se parece al de tu compañero. Tienes sprites, fondos, narrativa propios. (2) \emph{Coherencia}: las variables, condicionales y bucles tienen lógica clara, no están puestos al azar para cumplir requisitos. (3) \emph{Funcionamiento real}: el programa se ejecuta sin errores. Si los hay, los detectas en pruebas previas y los corriges.}
\milcpaso{4}{milcTurquesa}{\textbf{Errores típicos en la cosecha.} (1) \emph{Querer ser muy ambicioso} y no terminar a tiempo. Mejor un programa simple bien terminado que uno complejo a medias. (2) \emph{Empezar a programar sin plan} --- te pierdes en 30 minutos. Sigue tu plan de la Actividad 1. (3) \emph{Olvidar la documentación} --- es 10~\% de la rúbrica. (4) \emph{No probar el programa al final} --- llegas a la sustentación con errores. Prueba mínimo 5 veces. (5) \emph{Sustentar lectura, no demostración} --- en la sustentación, muestra el programa funcionando, no leas texto.}

\begin{drawbox}{milcTurquesa}{10 criterios + estructura del entregable}
\textbf{10 criterios:} plan · variables · condicionales · bucles · interactividad · diagrama · documentación · funciona · sustentación · personalización. \\[1mm]
\textbf{Entregable:} programa Scratch + diagrama cuaderno + documentación + sustentación 5 min.
\end{drawbox}

\begin{softbox}{milcMostaza}{milcGris}{Errores comunes y su corrección}
\textbf{Mitos sobre la cosecha del P2.} (1) \emph{``Mientras más complicado, mejor''} --- Falso. Mejor simple bien terminado que complejo incompleto. (2) \emph{``Si copio de un proyecto de Scratch que vi online, está bien''} --- Falso. Tu programa debe ser tuyo. Inspirarse sí; copiar no. (3) \emph{``La sustentación es solo presentación, no cuenta tanto''} --- Falso. Es 10~\% directo; las preguntas del profe pesan. (4) \emph{``El diagrama de flujo es opcional''} --- No. Es 10~\% de la rúbrica. (5) \emph{``Si funciona, la documentación es secundaria''} --- Falso. La documentación es donde demuestras qué entiendes, no solo qué hiciste.
\end{softbox}

\begin{infoband}{milcTurquesa}


\textbf{Lo que va al cuaderno (Actividad 2 · EXPLICA).} Título: «Actividad 2 --- Rúbrica + estructura del entregable». \textbf{Formato:} lista de los 10 criterios numerados + estructura del entregable. \textbf{Extensión:} 15 líneas. \textbf{Sabes que terminaste cuando} puedes autoevaluar cualquier programa Scratch con la rúbrica.
\end{infoband}

\puente{Con todo claro, ejecutas: programar + documentar + sustentar.}

\milcestacion{milcMagenta}{Estación 3 de 3 · Hacer}

\begin{softbox}{milcMagenta}{milcGris}{Cómo vas a construir tu producto}
\textbf{Actividad 3 · CREA --- Programa completo + diagrama + documentación + sustentación} (1.5 sesiones de 90 min + trabajo asíncrono durante la semana). Vas a ejecutar la cosecha completa. Es la prueba real del periodo.
\end{softbox}

\milcpaso{1}{milcMagenta}{\textbf{Paso 1 --- Programa el proyecto en Scratch (40-50 min).} Con tu plan de la Actividad 1 a la vista, abre Scratch (\texttt{scratch.mit.edu}) y construye el programa. Aplica los 5 elementos: 2 variables, 2 condicionales, 1 bucle, interactividad (preguntar), múltiples escenas/sprites. Personaliza con sprites y fondos propios. Prueba el programa al menos 3 veces durante la construcción para detectar errores temprano. Guarda con nombre \texttt{2026-cosecha-p2-[tu-nombre]}.}
\milcpaso{2}{milcMagenta}{\textbf{Paso 2 --- Dibuja el diagrama de flujo (15 min).} En una hoja del cuaderno, dibuja el diagrama de flujo aplicando S4: óvalo INICIO, rectángulos para procesos, rombos para condicionales (con SÍ/NO), paralelogramos para entradas/salidas, flechas direccionales, óvalo FIN. Si tu programa tiene bucles, dibújalos con la estructura de S7 (rombo con condición + flecha que regresa al rombo).}
\milcpaso{3}{milcMagenta}{\textbf{Paso 3 --- Documentación en el cuaderno (15 min).} Subtítulo: \textbf{``Documentación de mi programa''}. Incluye: \emph{(a) Descripción de 1 párrafo}: ¿qué hace tu programa, para qué sirve, quién es el usuario? \emph{(b) Lista de variables}: nombre + tipo + para qué sirve. \emph{(c) Capturas o bocetos de 3 momentos clave del programa} (escena inicial, una decisión importante, escena final). \emph{(d) Las 4 técnicas aplicadas}: 1 línea sobre cómo usaste descomposición, patrones, abstracción y algoritmo.}
\milcpaso{4}{milcMagenta}{\textbf{Paso 4 --- Autoevaluación con la rúbrica (5 min).} Antes de la sustentación, abre la rúbrica de 10 criterios. Marca cada uno con tu nota (1-5). Suma. Si algún criterio quedó bajo, ajusta si tienes tiempo. La autoevaluación honesta es parte del oficio adulto.}
\milcpaso{5}{milcMagenta}{\textbf{Paso 5 --- Sustentación de 5 minutos (en clase ante profe y grupo).} Estructura sugerida: \emph{(a) 30 segundos: presentar el tema del programa}. \emph{(b) 2 minutos: mostrar el programa funcionando (compartir pantalla en Teams o presencial)}. \emph{(c) 1 minuto: explicar las decisiones de diseño (por qué esa variable, ese condicional, ese bucle)}. \emph{(d) 30 segundos: contar qué te costó más y qué aprendiste}. \emph{(e) 1 minuto: preguntas del profe o compañeros}.}
\milcpaso{6}{milcMagenta}{\textbf{Paso 6 --- Cierre del periodo + firma final.} Después de la sustentación, escribe en el cuaderno tu reflexión de cierre del periodo: \emph{``En el periodo 2 aprendí: \_\_\_. Lo que más me sorprendió de la programación fue: \_\_\_. Si pudiera programar cualquier cosa, haría: \_\_\_''}. Firma con nombre y fecha. Cierra formalmente el periodo 2.}

\begin{drawbox}{milcMagenta}{Antes de sustentar}
\checkbox Mi programa Scratch funciona sin errores. \\[1mm]
\checkbox Tiene mínimo 2 variables, 2 condicionales, 1 bucle, interactividad. \\[1mm]
\checkbox El diagrama de flujo está dibujado en el cuaderno. \\[1mm]
\checkbox La documentación tiene descripción, variables, 3 capturas, 4 técnicas. \\[1mm]
\checkbox Hice autoevaluación con los 10 criterios de la rúbrica. \\[1mm]
\checkbox Estoy listo para sustentar 5 minutos con seguridad.
\end{drawbox}

\begin{softbox}{milcMagenta}{milcGris}{Plantilla de trabajo}
\textbf{Estructura.} \textit{Paso 1:} programa en Scratch (variables + condicionales + bucle + interactividad). \textit{Paso 2:} diagrama de flujo en cuaderno. \textit{Paso 3:} documentación. \textit{Paso 4:} autoevaluación con rúbrica. \textit{Paso 5:} sustentación 5 min. \textit{Paso 6:} reflexión de cierre.
\end{softbox}

\begin{infoband}{milcMagenta}


\textbf{Lo que va al cuaderno (Actividad 3 · CREA).} Título: «Actividad 3 --- Cosecha del P2: programa completo». \textbf{Formato:} diagrama + documentación + autoevaluación + reflexión. \textbf{Extensión:} 2 páginas. \textbf{Sabes que terminaste cuando} sustentaste el programa al profe y al grupo.
\end{infoband}

\puente{El producto es programa + diagrama + documentación + sustentación.}

\milcestacion{milcMostaza}{Producto de la sesión}
\textbf{Programa Scratch completo + diagrama + documentación + sustentación 5 min}.
\begin{softbox}{milcMostaza}{milcGris}{Criterios de calidad}
(1) El programa Scratch integra variables, condicionales, bucles, interactividad. (2) El diagrama de flujo está completo y correcto. (3) La documentación incluye descripción, variables, capturas y 4 técnicas. (4) Se hizo autoevaluación honesta con la rúbrica de 10 criterios. (5) La sustentación oral de 5 minutos fue clara y demostrativa.
\end{softbox}

\puente{Antes de sustentar, autoevalúa con la rúbrica de 10 criterios.}

\milcestacion{milcVino}{Evaluación liberadora · cinco dimensiones}

\begin{softbox}{milcVino}{milcGris}{Sentido de la evaluación}
Evaluar no es solo revisar una nota. Es reconocer qué aprendí, qué cambió en mí, qué emociones manejé, cómo me relaciono con la ciudad y la comunidad, y qué le dejaré a quien viene detrás.
\end{softbox}

\textbf{Mi reflexión liberadora en cinco dimensiones.} Completa cada frase iniciada:

\small
\begin{tabularx}{\textwidth}{>{\bfseries\RaggedRight\arraybackslash}p{3.4cm}Y>{\RaggedRight\arraybackslash}p{4.6cm}}
\rowcolor{milcVino}\color{white}Dimensión & \color{white}\bfseries Mi respuesta (frame starter) & \color{white}\bfseries Algo que descubrí\\
1. Personal & Lo que más cambió en mí fue \linead{6.5cm} & \linead{4.2cm}\\
2. Emocional & La emoción más fuerte fue \linead{2.5cm} y la regulé \linead{3cm} & \linead{4.2cm}\\
3. Ciudadana & En democracia esto importa porque \linead{6cm} & \linead{4.2cm}\\
4. Local (Cartago) & En mi barrio sería útil para \linead{6.5cm} & \linead{4.2cm}\\
5. Intergeneracional & A mi abuela/abuelo le diría \linead{6.5cm} & \linead{4.2cm}\\
\end{tabularx}
\normalsize

\begin{infoband}{milcVino}
\textbf{Para la próxima sustentación.} Vuelve a esta tabla en seis meses. Verás que las respuestas cambian — no porque lo aprendido cambie, sino porque tú cambias. La evaluación liberadora es una conversación contigo a través del tiempo.
\end{infoband}


\puente{Cerramos con 3 voces sobre el oficio cumplido y el camino hacia P3.}

\milcestacion{milcGradeDark}{Triángulo de pensamiento · cierre filosófico}

\begin{softbox}{milcVino}{milcGris}{Dussel · lente del nosotros}
\textit{El que termina su obra y la muestra a la comunidad cierra el ciclo. El que la oculta, la deja incompleta.}

{\footnotesize\itshape\color{milcNegro!70}--- formulación de esta guía, al modo de Enrique Dussel. No es una cita textual.}

Doña Sofía la tejedora del Valle no se quedaba con la mochila terminada en su casa: la \textbf{llevaba a la plaza}, la mostraba, recibía retroalimentación, transmitía el saber. \emph{Esa misma virtud aplica a tu cosecha hoy}. Tu programa de Scratch no es solo tu producto privado: es \emph{tu obra que muestras al grupo}. La sustentación cierra el ciclo del aprendizaje. \emph{Sin sustentar, el aprendizaje queda a medias}.

\textbf{¿Estoy mostrando mis obras (académicas y personales) o las guardo solo para mí?}
\end{softbox}
\linead{\linewidth}\\[1mm]
\linead{\linewidth}

\begin{softbox}{milcMostaza}{milcGris}{Estoicismo · Marco Aurelio · cuidado interior}
\textit{Termina lo que empezaste con la misma atención con que empezaste. La obra cumplida da paz; la inconclusa pesa siempre.}

{\footnotesize\itshape\color{milcNegro!70}--- formulación de esta guía, al modo de Marco Aurelio. No es una cita textual.}

Hace 10 sesiones empezaste el periodo \emph{sin saber programar}. Hoy entregas un programa con variables, condicionales, bucles, interactividad y sustentación. \textbf{Ese arco de 10 sesiones es ejemplo de oficio cumplido}. Marco Aurelio aplaudiría: \emph{empezar bien + sostener disciplina + cerrar bien}. La cosecha es la última disciplina, la que faltaba.

\textbf{¿Cuántas cosas empecé este año y dejé incompletas? ¿Qué pasaría si las terminara una a una?}
\end{softbox}
\linead{\linewidth}\\[1mm]
\linead{\linewidth}

\begin{softbox}{milcTurquesa}{milcGris}{Floridi · lente de la infoesfera}
\textit{Programar tu primer programa completo es rito de paso. Cambias de quien aprieta botones a quien crea con código. No hay vuelta atrás.}

{\footnotesize\itshape\color{milcNegro!70}--- formulación de esta guía, al modo de Luciano Floridi. No es una cita textual.}

\textbf{Hoy cambias de identidad}: antes eras \emph{usuario de programas}, ahora eres \emph{creador de programas}. La diferencia es enorme. \emph{Los usuarios consumen lo que otros producen; los creadores deciden qué se produce}. Tú hiciste un programa, no muy complejo, pero tuyo y completo. \emph{Esa transformación no se olvida}: tu yo de 25 años va a recordar que en grado 7 ya creabas con código. Esa identidad te acompañará a la universidad y al trabajo.

\textbf{Hoy soy creador digital. ¿En qué otras áreas de mi vida soy creador y no solo consumidor?}
\end{softbox}
\linead{\linewidth}\\[1mm]
\linead{\linewidth}

\begin{infoband}{milcNegro}
\textbf{Nota para el docente.} Las tres formulaciones son del autor de la guía, no citas textuales de los pensadores: recogen su lente para pensar el tema, no sus palabras. Si vas a citarlos en clase, remite a la obra original.
\end{infoband}

\milcestacion{milcVino}{Mi autoevaluación · marca una casilla por fila}

{\small
\setlength{\extrarowheight}{2.2pt}
\begin{tabularx}{\textwidth}{>{\RaggedRight\arraybackslash\bfseries}p{3.0cm}YYY}
\rowcolor{milcVino}\color{white}\bfseries Criterio & \color{white}\bfseries Lo logré & \color{white}\bfseries En proceso & \color{white}\bfseries Necesito apoyo\\[.6mm]
Comprensión del tema & \milcopcion{Explico las ideas con mis palabras.} & \milcopcion{Reconozco las ideas, falta precisión.} & \milcopcion{Necesito repasar lo básico.}\\[.8mm]
\rowcolor{milcGris}Pensamiento computacional & \milcopcion{Usé los pilares al armar mi producto.} & \milcopcion{Usé algunos pilares.} & \milcopcion{Necesito releer la estación 2.}\\[.8mm]
Aplicación y producto & \milcopcion{Mi producto cumple los criterios.} & \milcopcion{Está armado, con ajustes pendientes.} & \milcopcion{Aún está en construcción.}\\[.8mm]
\rowcolor{milcGris}Reflexión 5 dimensiones & \milcopcion{Respondí cada una con honestidad.} & \milcopcion{Respondí algunas con detalle.} & \milcopcion{Mis respuestas son muy generales.}\\[.8mm]
Triángulo de pensamiento & \milcopcion{Conecté las 3 voces con mi trabajo.} & \milcopcion{Conecté al menos una.} & \milcopcion{No las relaciono con mi proceso.}\\[.8mm]
\end{tabularx}}

\vspace{3mm}
\textbf{Hoy entendí que:}\\[1mm]
\linead{\linewidth}\\[1mm]
\linead{\linewidth}

\textbf{Mi compromiso para la próxima sesión es:}\\[1mm]
\linead{\linewidth}\\[1mm]
\linead{\linewidth}

\begin{softbox}{milcMostaza}{milcGris}{Cita de despedida · Marco Aurelio}
\textit{``El obstáculo en el camino se vuelve el camino. Lo que estorba al camino se vuelve camino.''}\\[1mm]
Cada error de hoy, bien observado, se convierte en pieza del camino que pisarás mañana.
\end{softbox}

\small
\begin{tabularx}{\textwidth}{>{\bfseries}p{3.6cm}Y}
\rowcolor{milcGradeDark!12}Nombre completo: & \linead{10cm}\\
Fecha: & \linead{4cm} \quad Curso/grupo: \linead{4cm}\\
Mi firma: & \linead{9cm}\\
\end{tabularx}
\normalsize

\begin{infoband}{milcGradeDark}
\textbf{Para la siguiente sesión.} Guarda tu programa en Scratch. La siguiente sesión es la \textbf{S11} (cosecha formal del periodo): autoevaluación con rúbrica visible, examen integrador del P2, y mirada al P3 que se viene: \emph{Inteligencia Artificial}. Hoy cerraste la programación básica; en P3 abres la puerta a la IA que ya está cambiando el mundo.
\end{infoband}

\end{document}
